\documentclass[11pt,a4paper]{article}
\usepackage{fontspec}
\usepackage[margin=1in]{geometry}
\usepackage[unicode,hidelinks]{hyperref}
\setlength{\parindent}{0pt}
\setlength{\parskip}{.6em}
\emergencystretch=3em
\newcommand{\volumeentry}[3]{\par\begingroup\parskip=.7em
  \noindent\href{Consciousness_Series_Complete_Collection.pdf\#page=#3}{\textbf{#1}\quad #2\nobreak\dotfill\nobreak\mbox{#3}}\par\endgroup}
\hypersetup{pdftitle={The Consciousness Series},pdfauthor={Sidong Liu}}
% COVER SUBTITLE: Realized access and first-person organization
\begin{document}
\begin{titlepage}
\centering
\vspace*{1.6cm}
{\Huge\bfseries The Consciousness Series\par}
\vspace{1cm}
{\Large Realized access and first-person organization\par}
\vfill
{\Large Complete Collected Volume\par}
\vspace{1.6cm}
{\Large Sidong Liu, PhD\par}
\vspace{.3cm}
{iBioStratix Ltd\par}
\href{mailto:sidongliu@hotmail.com}{\texttt{sidongliu@hotmail.com}}\par
\vspace{1cm}
{\large October 2026\par}
\end{titlepage}
\setcounter{page}{2}
\section*{Contents}
The numbers below refer to pages in this PDF. Each paper retains its own printed pagination, proofs, appendices and references.
\volumeentry{}{Introduction and reading map}{3}
\volumeentry{F-I}{Foundation I: From Source to Manifestation. Accessible Algebras from a Source, Its Noise and a Readable Seed}{4}
\volumeentry{F-II}{Foundation II: Shared Manifestation. Shared Structure, Public Calibration and Different Histories}{30}
\volumeentry{F-III}{Foundation III: The Perfuming Loop. Held Content, Learned Frames and Physical Return}{47}
\volumeentry{C-I}{Geometry Relative to a Cut: Selection, Generation and Modular Inseparability}{72}
\volumeentry{C-II}{Higher-Dimensional Consistency of Geometry and Modular Flow: Physical Source Identification and a Second-Order Extension}{96}
\volumeentry{C-III}{Finite Self-Models and the Explanatory Gap: What a Continuing Observer Can Know About Itself}{132}
\clearpage
\section*{Introduction}\label{collection:intro}
The six papers join Foundation's questions of realization, sharing and return to Consciousness's questions of geometry, perspective and self-description. Their common problem is how physical relations become available to a historically organized system, and what that availability explains.

The positive interfaces are specific. A selected hull can contain an internal standard Cut under the proper state and representation conditions. A shared conformal algebra with authorized operations can supply a common finite-energy test pulse. Its weak-mixture source can enter the physical identification and modular-consistency arguments. A finite history can select later access or protect held content, and an acquired record can enter a finite semiclassical mean-action source--response--reading return. These are not interchangeable definitions or a single universal apparatus.

The experience bridge is coordinated in Consciousness III rather than counted as several independent validations. It must pass through actual internal instruments and causal capacities, and then confront independent markers and alternatives. Brain Anchor and the four Imprint papers remain independent companion contributions: learned routing, damage and repair, source legibility across scales and the use of joint traces. They are not absorbed into Foundation III's finite return. The functional tests distinguish display-only explanations, but not every general metacognitive controller; independent reflective evidence remains necessary.

\section*{Reading map}
\begin{description}
\item[F-I] Source, noise and authorized seeds generate hulls of different types. Geometric inheritance requires the declared state, representation and internal standard Cut.
\item[F-II] Shared structure, shared state and correct public content are distinguished. Common finite-energy preparations provide a physical source-identification interface.
\item[F-III] Held content and finite physical return are realized with explicit resources. A history-selected frame and protected old quantum content are not the same code.
\item[C-I] Internal standard Cuts determine modular geometry in the stated class; the perspective argument distinguishes real access restrictions from recoding.
\item[C-II] Prepared source families and modular data test higher-dimensional geometry. Positive-energy Radon identifiability and common-source conditions keep the field-law claim conditional.
\item[C-III] Finite self-models have permission-dependent limits; one two-layer bridge connects internal geometric access to reflective organization and its independent empirical tests.
\end{description}
\end{document}
